\begin{textsample}{POS Dim 3 – ms – Score 35.00 – ms/ms\_inf\_e\_ef\_26.txt}  \label{ex:f3_pos_191}
7.3 Campos de Experiências Considerando os direitos acima elencados, a organização curricular da Educação \textbf{Infantil} na BNCC está estruturada em cinco campos de experiências e, a partir deles, foram definidos os objetivos de aprendizagem e desenvolvimento das \textbf{crianças}.
Essa forma de organização curricular requer a articulação dos saberes das \textbf{crianças}, associados às suas experiências, aos conhecimentos que fazem parte do patrimônio artístico, cultural, ambiental, científico e tecnológico, em consonância com as orientações das DCNEI ( 2009 ).
Considerando a ampliação desses saberes e conhecimentos e a concepção de currículo definida nas DCNEI ( 2009 ), os campos de experiências nos quais se organiza a BNCC são : O EU, O OUTRO E O NÓS É na interação com os \textbf{pares} e com \textbf{adultos} que as \textbf{crianças} vão constituindo um modo próprio de agir, sentir e pensar e vão descobrindo que existem outros modos de vida, pessoas diferentes, com outros pontos de vista.
Conforme vivem suas primeiras experiências sociais ( na família, na \textbf{instituição} escolar, na coletividade ), constroem \textbf{percepções} e questionamentos sobre si e sobre os outros, diferenciando -se e, simultaneamente, identificando -se como seres individuais e sociais.
Ao mesmo tempo que participam de relações sociais e de \textbf{cuidados} pessoais, as \textbf{crianças} constroem sua autonomia e senso de autocuidado, de reciprocidade e de interdependência com o meio.
Por sua vez, na Educação \textbf{Infantil}, é preciso criar oportunidades para que as \textbf{crianças} entrem em contato com outros grupos sociais e culturais, outros modos de vida, diferentes atitudes, técnicas e rituais de \textbf{cuidados} pessoais e do grupo, costumes, celebrações e narrativas.
Nessas experiências, elas podem ampliar o modo de perceber a si mesmas e ao outro, valorizar sua identidade, respeitar os outros e reconhecer as diferenças que nos constituem como \textbf{seres} humanos.
CORPO, \textbf{GESTOS} E MOVIMENTOS Com o corpo ( por meio dos sentidos, \textbf{gestos}, movimentos impulsivos ou intencionais, coordenados ou espontâneos ), as \textbf{crianças}, desde cedo, exploram o mundo, o espaço e os objetos do seu entorno, estabelecem relações, expressam -se, \textbf{brincam} e produzem conhecimentos sobre si, sobre o outro, sobre o universo social e cultural, tornando -se, progressivamente, conscientes dessa corporeidade.
Por meio das diferentes linguagens, como a música, a dança, o teatro, as \textbf{brincadeiras} de faz de conta, elas se comunicam e se expressam no entrelaçamento entre corpo, emoção e linguagem.
As \textbf{crianças} conhecem e reconhecem as sensações e funções de seu corpo e, com seus \textbf{gestos} e movimentos, identificam suas potencialidades e seus limites, desenvolvendo, ao mesmo tempo, a consciência sobre o que é seguro e o que pode ser um risco à sua integridade física.
Na Educação \textbf{Infantil}, o corpo das \textbf{crianças} ganha \textbf{centralidade}, pois ele é o \textbf{partícipe} privilegiado das práticas pedagógicas de \textbf{cuidado} físico, orientadas para a emancipação e a liberdade, e não para a submissão.
Assim, a \textbf{instituição} escolar precisa promover oportunidades ricas para que as \textbf{crianças} possam, sempre animadas pelo espírito \textbf{lúdico} e na interação com seus \textbf{pares}, explorar e vivenciar um amplo repertório de movimentos, \textbf{gestos}, olhares, \textbf{sons} e \textbf{mímicas} com o corpo, para descobrir variados modos de ocupação e uso do espaço com o corpo ( tais como sentar com apoio, rastejar, engatinhar, escorregar, caminhar \textbf{apoiando} -se em berços, mesas e cordas, saltar, escalar, equilibrar -se, correr, dar cambalhotas, alongar -se etc. ).
TRAÇOS, \textbf{SONS}, CORES E FORMAS Conviver com diferentes manifestações artísticas, culturais e científicas, locais e universais, no cotidiano da \textbf{instituição} escolar, possibilita às \textbf{crianças}, por meio de experiências diversificadas, vivenciar diversas formas de expressão e linguagens, como as artes visuais ( \textbf{pintura}, modelagem, colagem, fotografia etc. ), a música, o teatro, a dança e o audiovisual, entre outras.
Com base nessas experiências, elas se expressam por várias linguagens, criando suas próprias produções artísticas ou culturais, exercitando a autoria ( coletiva e individual ) com \textbf{sons}, traços, \textbf{gestos}, danças, \textbf{mímicas}, encenações, canções, desenhos, modelagens, manipulação de diversos materiais e de recursos tecnológicos.
Essas experiências contribuem para que, desde muito \textbf{pequenas}, as \textbf{crianças} desenvolvam senso estético e crítico, o conhecimento de si mesmas, dos outros e da realidade que as cerca.
Portanto, a Educação \textbf{Infantil} precisa promover a participação das \textbf{crianças} em tempos e espaços para a produção, manifestação e apreciação artística, de modo a favorecer o desenvolvimento da sensibilidade, da criatividade e da expressão pessoal das \textbf{crianças}, permitindo que se apropriem e reconfigurem, permanentemente, a cultura e potencializem suas singularidades, ao ampliar repertórios e interpretar suas experiências e vivências artísticas. \textbf{ESCUTA}, FALA, PENSAMENTO E IMAGINAÇÃO Desde o \textbf{nascimento}, as \textbf{crianças} participam de situações comunicativas cotidianas com as pessoas com as quais interagem.
As primeiras formas de interação do \textbf{bebê} são os movimentos do seu corpo, o olhar, a postura corporal, o sorriso, o choro e outros recursos vocais, que ganham sentido com a interpretação do outro.
Progressivamente, as \textbf{crianças} vão ampliando e enriquecendo seu vocabulário e demais recursos de expressão e de compreensão, apropriando -se da língua materna – que se torna, pouco a pouco, seu veículo privilegiado de interação.
Na Educação \textbf{Infantil}, é importante promover experiências nas quais as \textbf{crianças} possam falar e \textbf{ouvir}, potencializando sua participação na cultura oral, pois é na \textbf{escuta} de histórias, na participação em conversas, nas descrições, nas narrativas elaboradas individualmente ou em grupo e nas implicações com as múltiplas linguagens que a \textbf{criança} se constitui ativamente como sujeito singular e pertencente a um grupo social.
Desde cedo, a \textbf{criança} manifesta \textbf{curiosidade} com relação à cultura escrita : ao \textbf{ouvir} e acompanhar a leitura de textos, ao observar os muitos textos que circulam no contexto familiar, comunitário e escolar, ela vai construindo sua concepção de língua escrita, reconhecendo diferentes usos sociais da escrita, dos gêneros, suportes e portadores.
Na Educação \textbf{Infantil}, a \textbf{imersão} na cultura escrita deve partir do que as \textbf{crianças} conhecem e das \textbf{curiosidades} que deixam transparecer.
As experiências com a literatura \textbf{infantil}, propostas pelo educador, mediador entre os textos e as \textbf{crianças}, contribuem para o desenvolvimento do gosto pela leitura, do \textbf{estímulo} à imaginação e da ampliação do conhecimento de mundo.
Além disso, o contato com histórias, contos, fábulas, poemas, cordéis etc. propicia a familiaridade com \textbf{livros}, com diferentes gêneros literários, a diferenciação entre ilustrações e escrita, a aprendizagem da direção da escrita e as formas corretas de manipulação de \textbf{livros}.
Nesse convívio com textos escritos, as \textbf{crianças} vão construindo hipóteses sobre a escrita que se revelam, inicialmente, em rabiscos e garatujas e, à medida que vão conhecendo letras, em escritas espontâneas, não convencionais, mas já \textbf{indicativas} da compreensão da escrita como sistema de representação da língua.
ESPAÇOS, TEMPOS, QUANTIDADES, RELAÇÕES E TRANSFORMAÇÕES As \textbf{crianças} vivem inseridas em espaços e tempos de diferentes dimensões, em um mundo constituído de fenômenos naturais e socioculturais.
Desde muito \textbf{pequenas}, elas procuram se situar em diversos espaços ( rua, bairro, cidade etc. ) e tempos ( \textbf{dia} e noite ; hoje, ontem e amanhã etc. ). \textbf{Demonstram} também \textbf{curiosidade} sobre o mundo físico ( seu próprio corpo, os fenômenos atmosféricos, os animais, as plantas, as transformações da natureza, os diferentes tipos de materiais e as possibilidades de sua manipulação etc. ) e o mundo sociocultural ( as relações de parentesco e sociais entre as pessoas que conhece ; como vivem e em que trabalham essas pessoas ; quais suas tradições e seus costumes ; a diversidade entre elas etc. ).
Além disso, nessas experiências e em muitas outras, as \textbf{crianças} também se \textbf{deparam} com quantidades, dimensões, medidas, comparação de \textbf{pesos} e de comprimentos, avaliação de distâncias, reconhecimento de formas geométricas, conhecimento e reconhecimento de numerais cardinais e ordinais etc., que igualmente aguçam a \textbf{curiosidade}.
Portanto, a Educação \textbf{Infantil} precisa promover experiências nas quais as \textbf{crianças} possam fazer observações, \textbf{manipular} objetos, investigar e explorar seu entorno, levantar hipóteses e consultar fontes de informação para buscar respostas às suas \textbf{curiosidades} e indagações.
Assim, a \textbf{instituição} escolar está criando oportunidades para que as \textbf{crianças} ampliem seus conhecimentos do mundo físico e sociocultural e possam utilizá -los em seu cotidiano.
A organização do trabalho por Campos de Experiência procura favorecer as inter-relações entre os conhecimentos, nas quais as vivências das \textbf{crianças} não podem ser vistas de forma fragmentada, mesmo porque na Educação \textbf{Infantil} os tempos são diferenciados.
Nesse sentido, faz -se necessário que essas vivências sejam propiciadas pelo professor nas \textbf{instituições}.
Sob essa vertente cabe ressaltar : > Os saberes e conhecimentos prévios do professor, sua formação científica, artística, tecnológica, ambiental, cultural lhe possibilita enriquecer ou ampliar o currículo vivido pelas \textbf{crianças} no cotidiano da \textbf{creche} e da \textbf{pré-escola}. > ( BARBOSA ; RICHTER, 2015, p.196 ) Assim, nas propostas pedagógicas deverão ser estabelecidas a identidade da \textbf{instituição} e as escolhas pedagógicas nas quais os saberes e conhecimentos de diferentes naturezas que compõem os Campos de Experiências e suas subdivisões internas possibilitem aprendizagem e desenvolvimento a todas as \textbf{crianças}.

% matched lemmas: adulto, apoiar, bebê, brincadeira, brincar, centralidade, creche, criança, cuidado, curiosidade, demonstrar, deparar, dia, escuta, estímulo, gesto, imersão, indicativo, infantil, instituição, livro, lúdico, manipular, mímico, nascimento, ouvir, par, partícipe, pequeno, percepção, peso, pintura, pré-escola, seres, som
\end{textsample}
